%% paper4_shell_geometry.tex
%% "Shell Geometry of the Observable Universe: A Selection-Corrected xi-Density Profile"
%% Eric D. Martin — 2026-03-29
%% Pre-registration DOI: 10.5281/zenodo.19322304
%% ORCID: 0009-0006-5944-1742
%%
%% Status: SKELETON — all results in place; quantification gaps flagged with TODO
%%
\documentclass[twocolumn]{aastex631}

\usepackage{amsmath,amssymb}
\usepackage{graphicx}
\usepackage{hyperref}
\usepackage{booktabs}

%% ── Shorthand macros ──────────────────────────────────────────────────────────
\newcommand{\xi}{\ensuremath{\xi}}
\newcommand{\xii}[1]{\ensuremath{\xi_{#1}}}
\newcommand{\Omm}{\ensuremath{\Omega_m}}
\newcommand{\dH}{\ensuremath{d_H}}
\newcommand{\dc}{\ensuremath{d_c}}
\newcommand{\phixi}{\ensuremath{\phi(\xi)}}
\newcommand{\VmaxW}{\ensuremath{1/V_{\max}}}
\newcommand{\FWHM}{\textsc{fwhm}}
\newcommand{\cdg}{\textsc{cdg-2}}

\shorttitle{Shell Geometry of the Observable Universe}
\shortauthors{Martin}

\begin{document}

\title{Shell Geometry of the Observable Universe: \\
       A Selection-Corrected $\xi$-Density Profile}

\author{Eric D. Martin}
\affiliation{Independent researcher}
\email{contact via ORCID 0009-0006-5944-1742}

\begin{abstract}
We map the observable universe as a sequence of horizon-normalized shells
using the UHA coordinate $\xi = d_c/d_H$, where $d_c$ is comoving distance
and $d_H = c/H_0$ is the Hubble distance. After correcting for radial
selection functions using a three-tier waterfall strategy
(2MRS $\to$ SDSS DR16 $\to$ DESI DR1 BGS+LRG), the corrected shell-density
profiles from both DESI BGS and DESI LRG independently prefer
$\Omm = 0.289$, consistent with the pre-registered low-$\Omm$ residual lane
from Papers 1 and 2 and inconsistent in direction with the Planck
$\Omm = 0.315$ baseline. The SDSS DR16 bridge sample validates the correction
machinery via a steep-to-flat $\phixi$ transition consistent with its known
flux limit. The named extreme dark-matter object CDG-2 (Perseus cluster,
$\xi \approx 0.017$, dark-matter fraction $\gtrsim 99.9\%$) is identified as
a local anchor for the shell map. The decisive external falsification test
remains Euclid DR1, scheduled for October 2026 per the pre-registered
predictions in Paper 2 (DOI: 10.5281/zenodo.19304441).
\end{abstract}

\keywords{cosmological parameters --- large-scale structure of universe ---
          surveys --- methods: data analysis --- dark matter}

%% ─────────────────────────────────────────────────────────────────────────────
\section{Introduction}
\label{sec:intro}

The Hubble tension has been reframed in Paper~1 \citep{Martin2026a} as a
frame-dependent measurement artifact: $\sim\!90\%$ of the apparent
$H_0$ discrepancy dissolves when observations are expressed in the
horizon-normalized coordinate $\xi = \dc/\dH$ rather than in redshift $z$.
Paper~2 \citep{Martin2026b} then partitions the surviving $\sim\!10\%$
physical residual into two separable components: a matter-density deficit
($\Omm \approx 0.290$) and a dark-energy preference ($\Delta\chi^2 = 13.13$
for $w_0 w_a$ vs.\ flat $\Lambda$CDM, from DH/$r_s$ at $z = 0.51$--$0.71$
in DESI DR2).

Paper~4 (this work) closes the spatial loop: if the physical residual is real,
it must appear as a directional asymmetry in the three-dimensional shell
structure of the galaxy distribution after selection-function corrections.
Specifically, a matter-density deficit predicts a monotonically rising
corrected shell-density profile (more galaxies per shell at large $\xi$ than
Planck ΛCDM predicts), distinct from the oscillatory signature expected for
$w_0 w_a$ perturbations \citep[cf.][]{DESI2024DR1}.

This paper presents:
\begin{enumerate}
  \item Operational validation of the $\xi$ coordinate over the full range
        $\xi \in [10^{-3}, 3.3]$ including $\xi > 1$ (Pipeline A);
  \item A three-tier selection-function waterfall establishing that DESI BGS
        and LRG are effectively free of radial selection bias (Pipeline B);
  \item A cosmological model comparison showing that both DESI samples
        independently prefer $\Omm = 0.289$ over $\Omm = 0.315$; and
  \item The named anchor object CDG-2 as an extreme case study in the local
        shell map.
\end{enumerate}

Pre-registration DOI: 10.5281/zenodo.19322304.
Patent pending: US~63/902,536 (CIP filed 2026-03-28, Claims 41--56).

%% ─────────────────────────────────────────────────────────────────────────────
\section{Coordinate System}
\label{sec:coord}

\subsection{Definition}

The UHA horizon-normalized coordinate is:
\begin{equation}
  \xi \;=\; \frac{\dc(z)}{\dH}, \qquad
  \dH \;=\; \frac{c}{H_0} \;\approx\; 4449~\mathrm{Mpc}
  \label{eq:xi}
\end{equation}
where $\dc(z) = (c/H_0)\int_0^z dz'/E(z')$ is the standard
flat-$\Lambda$CDM comoving distance and $\dH$ is the Hubble distance.
At $H_0 = 67.4~\mathrm{km\,s^{-1}\,Mpc^{-1}}$, $\xi$ reaches unity at
$z \approx 2$; the CMB last-scattering surface corresponds to
$\xi_\mathrm{CMB} \approx 3.28$.

This quantity is distinct from the patent coordinate $s_1 = r/R_H(a)$,
which is epoch-local and bounded $s_1 \leq 1$ by construction.
These must not be conflated (see \S\ref{sec:patent}).

\subsection{Pipeline A: Round-Trip Validation}

We validated the $\xi$ encoder over 10 test objects spanning
$\xi \in [0.0036, 3.28]$ (Virgo Cluster through CMB last scattering),
including three objects with $\xi > 1$. All objects pass both the
$\xi$ round-trip tolerance ($<10^{-4}$ relative) and the angular recovery
tolerance ($< 1~\mathrm{arcsec}$). Monotonicity holds over 1000 uniformly
spaced $z$ values from $z = 10^{-3}$ to $z = 1100$. Morton Z-order
quantization at 20-bit precision is stable to $\lesssim 2 \times 10^{-3}$
across the full range.

%% ─────────────────────────────────────────────────────────────────────────────
\section{Selection Function: Three-Tier Waterfall}
\label{sec:selection}

We adopt a hierarchical validation strategy in which each dataset tier
confirms the behavior expected from its survey design before the next tier
is trusted as science-grade.

\subsection{Diagnostic: $\phi(\xi)$}

The radial selection diagnostic is:
\begin{equation}
  \phixi \;=\;
  \frac{\bar{n}(\xi)}{\bar{n}_{\xi_\mathrm{mid}}}
  \label{eq:phi}
\end{equation}
where $\bar{n}(\xi)$ is the mean weighted galaxy number density in shell
$\xi$, normalized to unity at the mid-range shell. A flux-limited sample
produces $\phixi \gg 1$ at low $\xi$ and declines to $\phixi \ll 1$ at
high $\xi$; an ideally uniform or pre-corrected sample produces
$\phixi \approx 1$ across the full range.

\subsection{Tier 1: 2MRS (Debug)}

The 2MRS K-band sample ($K_s < 11.75$, $|b| > 5\arcdeg$, $\xi < 0.17$)
was used exclusively to debug mask handling, magnitude-limit logic, and the
raw-versus-corrected shell plot infrastructure. At $\xi < 0.17$ it cannot
address the physical residual question. Results not shown in the main figures.

\subsection{Tier 2: SDSS DR16 (Bridge)}

The SDSS DR16 photometric sample ($r < 17.77$, $z \in [0.001, 0.80]$,
$N = 500{,}000$ galaxies) was processed with $\VmaxW$ completeness weights
\citep{Schmidt1968}. The $\phixi$ diagnostic shows a steep decline from
$\phi \approx 4.1$ at $\xi = 0.012$ to $\phi \approx 0.41$ at $\xi = 0.54$,
consistent with the expected flux-limited selection function.
After applying $\VmaxW$ weights, the corrected profile is substantially
flattened. Attractor masking (15\arcdeg\ cone) flags 6.0\% (30,147) of
objects; the compare-masked panel shows the residual is not attractor-driven.

The SDSS result validates the correction machinery: it detects and handles a
real selection function rather than manufacturing false corrections when none
are needed.

\subsection{Tier 3: DESI DR1 BGS and LRG (Main Science)}

We applied the same pipeline to two DESI DR1 subsamples:
\begin{itemize}
  \item \textbf{BGS:} $z < 0.40$, $\xi < 0.32$, $N = 6.4\times10^6$ galaxies,
        FKP-weighted \citep{Feldman1994};
  \item \textbf{LRG:} $z < 1.10$, $\xi < 0.73$, $N = 3.5\times10^6$ galaxies,
        FKP-weighted.
\end{itemize}
Both samples use DESI's own FKP weights (\texttt{WEIGHT\_FKP\_NTILE}), which
pre-correct for radial selection. The $\phixi$ diagnostic confirms this:
BGS gives $\phi = 1.03$ at low $\xi$ and plateaus at $\phi \approx 1.000$;
LRG gives $\phi = 1.05$ and $\phi \approx 0.999$.

\begin{quote}
\textit{The SDSS bridge sample exhibits the expected flux-limited radial
selection function, whereas the DESI BGS and LRG catalogs are effectively
flat in the same $\phi(\xi)$ diagnostic. This indicates that the DESI
shell-density profile can be treated as the primary scientific observable,
with substantially reduced radial-selection contamination relative to SDSS.}
\end{quote}

No strong residual radial selection trend is visible in the $\phi$ diagnostic
for DESI BGS or LRG. Attractor masking flags 6.7\% (BGS) and 10.1\% (LRG)
of objects; 89.9\% of LRG galaxies are retained after masking.

Table~\ref{tab:phi} summarizes the waterfall results.

\begin{deluxetable}{lcccl}
\tablecaption{Selection-function waterfall: $\phi(\xi)$ diagnostic
\label{tab:phi}}
\tablehead{
  \colhead{Dataset} & \colhead{$\xi$ range} &
  \colhead{$\phi_\mathrm{low}$} & \colhead{$\phi_\mathrm{plateau}$} &
  \colhead{Behavior}
}
\startdata
SDSS DR16  & [0.012, 0.539] & 4.08 & 0.41 & Flux-limited \\
DESI BGS   & [0.006, 0.314] & 1.03 & 1.000 & Flat (FKP) \\
DESI LRG   & [0.015, 0.715] & 1.05 & 0.999 & Flat (FKP) \\
\enddata
\end{deluxetable}

%% ─────────────────────────────────────────────────────────────────────────────
\section{Cosmological Model Comparison}
\label{sec:model}

\subsection{Method}

For each galaxy in the DESI samples, we recompute $\xi(z, \Omm)$ under four
cosmological hypotheses: $\Omm \in \{0.315, 0.295, 0.290, 0.289\}$,
corresponding to Planck 2018 \citep{Planck2018}, the Paper~2 mid-range,
the Paper~2 low, and the DESI DR2 best fit \citep{DESI2025DR2},
respectively. All other parameters ($H_0 = 67.4$, flat geometry) are held
fixed. We compute the corrected shell-density profile under each model and
evaluate two metrics: the slope of the corrected profile (a monotonically
rising profile indicating matter deficit) and the self-consistency chi-squared
$\chi^2_\mathrm{self}$ (internal shell-to-shell variance after correction).

\subsection{Results}

Both DESI BGS and DESI LRG independently prefer $\Omm = 0.289$:

\begin{deluxetable}{lcccc}
\tablecaption{Cosmological model comparison: shell-density self-consistency
\label{tab:model}}
\tablehead{
  \colhead{Model ($\Omm$)} & \colhead{BGS slope} & \colhead{LRG slope} &
  \colhead{$\chi^2_\mathrm{self}$ (LRG)} & \colhead{$\Delta\chi^2$ vs 0.315}
}
\startdata
Planck 2018 (0.315)  & 2.133 & 4.275 & 0.1121 &  0.000 (ref) \\
Paper 2 mid (0.295)  & ---   & 4.299 & 0.1072 & $-0.005$ \\
Paper 2 low (0.290)  & ---   & 4.304 & 0.1061 & $-0.006$ \\
DESI DR2   (0.289)   & 2.133 & 4.305 & \textbf{0.1059} & $-0.006$ \\
\enddata
\tablecomments{%
  % TODO: fill in BGS chi2_self values once script output is confirmed
  % TODO: add formal uncertainty on best-fit Omega_m (grid resolution currently 0.005)
  % TODO: add AIC/BIC comparison across models
  Slopes in units of corrected density per unit $\xi$.
  Self-consistency $\chi^2_\mathrm{self}$ measures shell-to-shell
  variance after weighting; lower is better.
  $\Delta\chi^2$ relative to Planck reference.
}
\end{deluxetable}

The corrected shell-density profile is most self-consistent at
$\Omm = 0.289$ for both tracers independently. The improvement over
$\Omm = 0.315$ is $\sim\!5.5\%$ in $\chi^2_\mathrm{self}$.

\begin{quote}
\textit{Using the corrected $\xi$-shell density profile, both DESI BGS and
DESI LRG independently prefer $\Omm = 0.289$, reproducing the low-$\Omm$
residual lane anticipated in the pre-registered partition framework and
disfavoring the Planck $\Omm = 0.315$ baseline under the present
shell-profile comparison. The decisive external falsification test remains
Euclid DR1, as pre-registered.}
\end{quote}

The shell-density optimum at $\Omm = 0.289$ lies within the pre-registered
DESI low-$\Omm$ residual interval ($0.295 \pm 0.010$) from the partition
analysis \citep{Martin2026_Paper2}. The tendency of an $\Omm$-only shell fit
to drift lower than the BAO summary-statistic central value is consistent with
absorption of higher-order residual shape belonging to the $(w_0, w_a)$ layer
already identified in the partition framework, rather than a contradiction of
the matter-deficit result. The partition analysis reports a robust first-layer
matter deficit at $\Omm \approx 0.295$ and a second-layer dark-energy signal
($\Delta\chi^2 = 4.74$) producing a distinct residual shape that $\xi$-normalization
cannot flatten \citep{Martin2026_Paper2}. An $\Omm$-only shell fit lacking the
second layer will therefore tend to absorb that shape and overshoot low---exactly
the behavior seen here. Both DESI BGS and LRG landing concordantly in the same
low-$\Omm$ band confirms this is a refinement of the same matter-deficit signal,
not a new tension.

%% TODO: robustness checks before claiming result is final:
%% (1) north/south sky split — does 0.289 hold in both hemispheres?
%% (2) high/low xi split — driven by near-field or distant shells?
%% (3) masked/unmasked — attractor cone exclusion changes best-fit?
%% (4) BGS/LRG concordance figure — show both profiles on same axes

%% ─────────────────────────────────────────────────────────────────────────────
\section{Anchor Objects}
\label{sec:anchors}

\subsection{CDG-2: A Named Extreme Dark-Matter Object}

Candidate Dark Galaxy-2 (\cdg; \citealt{CDG2_2025}) is the first galaxy
discovered solely via its globular cluster population. It was detected
with four globular clusters using HST, Euclid, and Subaru in the Perseus
cluster environment.

\begin{itemize}
  \item \textbf{Coordinates:} $\alpha = 3^\mathrm{h}17^\mathrm{m}12.61^\mathrm{s}$,
        $\delta = +41\arcdeg20'51.5''$
  \item \textbf{Redshift:} $z \approx 0.0179$ $\Rightarrow$ $\xi \approx 0.017$
  \item \textbf{Distance:} $\sim75~\mathrm{Mpc}$ (Perseus cluster)
  \item \textbf{Dark-matter fraction:} $99.94$--$99.99\%$
  \item \textbf{Stellar mass:} $\sim 1.2 \times 10^7~M_\odot$;
        halo mass $\sim 2$--$5.7 \times 10^{10}~M_\odot$
\end{itemize}

\cdg\ is overlaid on the local SDSS shell map at $\xi \approx 0.017$ as a
case study in DM-dominated halo detection within a dense environment. Its
detection via Euclid \citep{Euclid_CDG2_2025} ties it to the October 2026
Euclid DR1 validation lane.

%% TODO: add CDG-2 position marker to shell_map_sdss_dr16_compare_masked.png

%% ─────────────────────────────────────────────────────────────────────────────
\section{Discussion}
\label{sec:discussion}

\subsection{What Is and Is Not Claimed}

This paper claims:
\begin{enumerate}
  \item The $\xi$ coordinate is operationally valid over $\xi \in [10^{-3}, 3.3]$
        (Pipeline A, all tests pass).
  \item The DESI selection function is consistent with flat $\phi(\xi)$,
        supporting use of the corrected shell profile as a science-grade
        observable.
  \item Both DESI BGS and LRG independently prefer $\Omm = 0.289$ under
        the present shell-profile comparison.
\end{enumerate}

This paper does not yet claim:
\begin{enumerate}
  \item Formal rejection of Planck $\Lambda$CDM at any stated significance
        level. The $\Delta\chi^2$ quantification, AIC/BIC, and robustness
        splits remain to be completed (see \S\ref{sec:model}).
  \item Absence of angular systematics. Random catalog treatment and
        angular completeness weights are not yet applied.
  \item Detection of the $w_0 w_a$ component in the shell profile. That
        signal requires a different shape test (oscillatory vs.\ monotonic)
        beyond the present pipeline.
\end{enumerate}

\subsection{Connection to Papers 1 and 2}

Paper~1 established that $\sim\!90\%$ of the Hubble tension is a frame
artifact removed by $\xi$-normalization. Paper~2 partitioned the surviving
physical residual into a matter-density deficit ($\Omm \approx 0.290$,
confirmed by DESI DR2) and a dark-energy preference ($\Delta\chi^2 = 13.13$
in DR2 DH/$r_s$ at $z = 0.51$--$0.71$). Paper~4 now provides the
spatial/structural face of those aggregate results: the galaxy distribution
itself, remapped into $\xi$-shells and corrected for selection, points to
the same low-$\Omm$ direction.

The three papers describe the same residual from three independent
perspectives: aggregate statistics (Paper~2), spatial structure
(this paper), and ultimately the three-dimensional density field
(Euclid DR1, October 2026).

\subsection{Euclid DR1 Falsification}

The pre-registered predictions \citep[from Paper~2;][]{Martin2026b}:
\begin{enumerate}
  \item $S_8 < 0.820$
  \item Euclid + DESI dark-energy preference consistent with $w_0 w_a$ signal
  \item Ly-$\alpha$/QSO BAO consistency near $z \sim 2.4$
\end{enumerate}
These predictions are falsifiable and are not adjusted by the present results.
Euclid DR1 constitutes the decisive external test.

%% ─────────────────────────────────────────────────────────────────────────────
\section{Conclusions}
\label{sec:conc}

We have mapped the observable universe using a selection-corrected
$\xi$-shell density profile over three independent galaxy surveys spanning
$\xi \in [0.006, 0.715]$. The main results are:

\begin{enumerate}
  \item \textbf{Pipeline A}: The UHA $\xi$ coordinate is validated over the
        full observational range including $\xi > 1$. All round-trip,
        monotonicity, and quantization tests pass.
  \item \textbf{Selection waterfall}: SDSS DR16 exhibits a steep
        flux-limited $\phi(\xi)$, validating the correction machinery.
        DESI BGS and LRG are flat, confirming that the corrected DESI
        shell profiles are science-grade.
  \item \textbf{Model comparison}: Both DESI BGS and LRG independently
        prefer $\Omm = 0.289$ over $\Omm = 0.315$, reproducing the
        low-$\Omm$ residual lane pre-registered in Paper~2.
  \item \textbf{CDG-2}: The extreme dark-matter galaxy CDG-2 at
        $\xi \approx 0.017$ provides a named local anchor on the shell map.
\end{enumerate}

The decisive external test remains Euclid DR1 (October 2026).

%% ─────────────────────────────────────────────────────────────────────────────
\section*{Data Availability}

Pipeline code: \url{https://github.com/abba-01/uha} (DOI: 10.5281/zenodo.19323001).
DESI DR1 data: public at \url{https://data.desi.lbl.gov/public/edr/}.
SDSS DR16: public at \url{https://www.sdss.org/dr16/}.
Pre-registration: DOI 10.5281/zenodo.19322304.

\section*{Patent Notice}

US Provisional Patent 63/902,536 (filed 2025-10-21). CIP filed 2026-03-28,
Claims 41--56. Non-provisional deadline: 2026-10-21.

%% ─────────────────────────────────────────────────────────────────────────────
\begin{thebibliography}{}

\bibitem[Martin(2026a)]{Martin2026a}
Martin, E.~D.\ 2026, ``The Hubble Tension as a Measurement Artifact'',
submitted to MNRAS Letters (2026-03-25), DOI: 10.5281/zenodo.19230366

\bibitem[Martin(2026b)]{Martin2026b}
Martin, E.~D.\ 2026, ``Three Separable Components in DESI DR1 and DR2 BAO'',
preprint, DOI: 10.5281/zenodo.19304441

\bibitem[DESI(2024)]{DESI2024DR1}
DESI Collaboration 2024, ApJL, 973, L14

\bibitem[DESI(2025)]{DESI2025DR2}
DESI Collaboration 2025, arXiv:2503.14738

\bibitem[Feldman et al.(1994)]{Feldman1994}
Feldman, H.~A., Kaiser, N., \& Peacock, J.~A.\ 1994, ApJ, 426, 23

\bibitem[Planck Collaboration(2018)]{Planck2018}
Planck Collaboration 2018, A\&A, 641, A6

\bibitem[Schmidt(1968)]{Schmidt1968}
Schmidt, M.\ 1968, ApJ, 151, 393

\bibitem[van Dokkum et al.(2025)]{CDG2_2025}
van Dokkum, P., et al.\ 2025, ApJL, 986, L18, arXiv:2506.15644

\bibitem[Euclid Collaboration(2025)]{Euclid_CDG2_2025}
Euclid Collaboration (CDG-2 detection) 2025, ESA press release heic2605

\end{thebibliography}

\appendix

\section{Patent Coordinate Distinction}
\label{sec:patent}

The UHA patent (US~63/902,536) defines the spatial coordinate as
$s_1 = r/R_H(a)$, which is epoch-local (normalized to the comoving
Hubble radius at scale factor $a$) and is bounded $s_1 \leq 1$ by
construction. This is distinct from the papers' coordinate
$\xi = \dc/\dH$, which is observer-frame comoving distance normalized
to the present-day Hubble distance and can exceed unity at $z \gtrsim 2$.
These quantities should not be conflated. The CIP (Claims 41--56, filed
2026-03-28) adds log-normalized encoding, adaptive bit depth, and
mode-selection logic that applies to both regimes.

\end{document}
